\section{Explanation Ground-Truth Generation}
\label{sec:ground-truth-generation}
%
% PURPOSE OF THIS CHAPTER
% Define every explanation reference produced from the known generator, its
% mathematical object, its indices, and the class of explanation that may be
% compared with it.  Keep process explanations distinct from explanations of a
% fitted forecasting model.
%
% \subsection{Ground-truth contract and shared indices}
% - Define instance n, target k, horizon h, source series j, and lag ell once.
% - Specify shapes, dtypes, missing/not-applicable values, and metadata.
% - State whether a reference describes one-step structure, an h-step unrolled
%   mechanism, a local realized effect, or an attribution under a baseline/game.
%
% \subsection{Structural source-lag ground truth}
% - Explain extraction of the binary source-lag mask directly from the AST.
% - Explain how recursively generated intermediate forecasts propagate dependencies
%   to each horizon instead of copying the one-step mask to every horizon.
% - Include small hand-derived one-step and multi-step examples.
%
% \subsection{Signed coefficients and realized local terms}
% - Define when an AST admits an identifiable coefficient reference.
% - Separate the signed parameter from the realized coefficient-times-input term.
% - Mark these references not applicable for operators without the required
%   coefficient interpretation.
%
% \subsection{Gradient-based references}
% - Define the derivative of the h-step generator output with respect to every
%   input variable-lag coordinate.
% - Explain that recursive differentiation automatically includes paths through
%   earlier forecast steps.
% - Separate raw gradient (local sensitivity) from baseline-relative
%   gradient-times-input (a contribution heuristic).
% - Document differentiability restrictions and finite-difference validation.
%
% \subsection{Generator-level Shapley references}
% - Define players, coalitions, the generator as the value-producing function,
%   the background distribution, and the rule used when a feature is absent.
% - Explain exact enumeration for small games and permutation estimation for
%   larger games.
% - Describe dependency-aware resampling and grouped/Owen-value extensions only
%   if they are implemented in the final benchmark.
% - Retain signed local attributions per forecast horizon before any absolute or
%   global aggregation.
%
% \subsection{Validation of reference explanations}
% - List hand-derived formulas, efficiency/additivity checks when applicable,
%   autodiff versus finite differences, and exact versus sampled Shapley checks.
% - Explain how seeds and estimator uncertainty are recorded.
% - End with a compatibility table mapping each ground-truth family to valid
%   comparison metrics and model/explainer outputs.
